% !TEX root = ../main.tex
\section{Methods}

\hypertarget{case-study-the-umpqua-national-forest-forplan-model}{%
\subsection{Case study: the Umpqua National Forest FORPLAN model}\label{case-study-the-umpqua-national-forest-forplan-model}}

The case study reproduces the Umpqua National Forest FORPLAN planning model \citep{usda1987umpqua}, reconstructed in the open-source \texttt{ws3} wood-supply model \citep{paradis2026ws3}. The forest is stratified into four ecoclasses -- CH-CW, CD-CP, CR-CF, and CM-CE, dominated respectively by Douglas-fir, hemlock, fir, and hemlock types -- arrayed in decreasing growth potential. Seven silvicultural prescriptions are available, each with ecoclass-specific rotation-age ranges. The planning horizon is 150 years, divided into 15 periods of 10 years. The objective is maximization of net present value at a 4\% discount rate (varied in the experiments), with delivered-log prices escalating at 1\% per year for the first 50 years and constant thereafter. Volumes are expressed in thousands of cubic feet (MCF). The harvest-flow policy is the thesis's bounded-deviation \emph{sequential-flow} constraint of equation~\eqref{eq:openloop}, tying each period's total harvest volume to the previous period's; the specific policy levels are the thesis's six harvest-flow constraint sets (Section~\ref{experiment-design}). The initial forest is specified as one of a set of landbase conditions; the landbase of central interest (landbase 1) is an all-mature forest, the extreme of the disequilibrium structures the thesis studies.

\hypertarget{case-study-data}{%
\subsection{Case-study data}\label{case-study-data}}

The case study is built from the real Umpqua FORPLAN data, obtained from the public-domain Umpqua National Forest planning documents on HathiTrust. The per-age yield curves -- the volume removed per acre by age for each ecoclass, prescription, and management intensity -- come from the 1990 Final Environmental Impact Statement (FEIS) Appendix B \citep{usda1990umpqua} (the FORPLAN analysis volume: managed yield tables generated with the DFSIM Douglas-fir simulator, with operational-falldown adjustments, and the mountain-hemlock curve from Johnson's mountain-hemlock site equations). The economics come from the FEIS's timber-value derivations: stumpage, logging, and manufacturing costs by species (its Table B-65), the Douglas-fir price-diameter and pond-value relationships (Table B-66), and the per-ecoclass access (road) costs that make the high-elevation mountain-hemlock (CM-CE) prescriptions negatively valued. The yield-curve shapes are consistent with the LRMP's documented 95\% CMAI culmination ages by ecoclass (its Table IV-3: 85 years for CH-CW, 120 for CD-CP, 115 for CR-CF, 175 for CM-CE), and the FEIS Stage II financial analysis corroborates the thesis's anchor values (its mature-timber NPV range of \$1,800--\$7,700 per acre brackets the thesis's Table 5.4 mature-type NPVs; its managed-prescription range of \$27--\$290 per acre and the negative mountain-hemlock values match the thesis's Table 5.3 ordering and signs).

All source documents were obtained and machine-read from HathiTrust (public domain): the 1990 LRMP, the 1990 FEIS main volume and its Appendix B (the FORPLAN analysis volume), and the 1987 Draft EIS \citep{usda1987umpqua} that Daugherty used directly (the thesis's exact data vintage). The 1987 DEIS and 1990 FEIS yield tables agree, so the assembled case study uses Daugherty's exact data. The case study reproduces the thesis's printed anchor tables (Tables 5.3 and 5.4): the five mature vegetation types' model PNVs match their Table 5.4 anchors exactly (their standing volumes are set so that volume $\times$ the FEIS net delivered-log price equals the anchor), and the managed prescriptions match Table 5.3 in sign and broad rotation range (the exact optimal-rotation ordering is only partially reproduced, as the FEIS prescription mapping simplifies some treatments; Section 5.5). Because the mature volumes are derived from the Table 5.4 anchors, we additionally cross-check them against an independent source -- the FEIS standing-volume-by-age curves at each mature type's age -- which agree in order of magnitude. All data are typed, provenance-stamped records in the repository; the extraction of the FEIS Appendix B tables is a one-time, documented, reproducible step, and the package does not depend on the scanned source at runtime.

\hypertarget{the-sequential-replanning-simulator}{%
\subsection{The sequential-replanning simulator}\label{the-sequential-replanning-simulator}}

Dynamic inconsistency cannot be observed in a single open-loop solve; it is defined by what happens when the plan meets sequential re-optimization. Following the thesis's iterative LP simulation, we implement a sequential-replanning simulator:

\begin{enumerate}
\def\labelenumi{\arabic{enumi}.}
\tightlist
\item
  Solve the open-loop LP from the current forest state over the horizon.
\item
  Apply the current period's decision from that solution.
\item
  Advance the forest state by one period: the realized area-by-age distribution is extracted from the model and becomes the initial condition of the next subproblem.
\item
  Re-solve the open-loop LP from the realized state over the (rolling) horizon.
\item
  Repeat to the end of the simulation, recording both the plan projected at each step and the decisions actually taken.
\end{enumerate}

Two modeling decisions bear on the interpretation and are worth stating explicitly. First, each replan re-solves the harvest-scheduling problem with a \emph{fresh} harvest-flow constraint anchored to that subproblem's own trajectory, not carried from the prior period's realized harvest. This is the faithful reading of FORPLAN practice -- each decadal re-solve is a new plan with its own flow constraint -- and it is what allows the promised non-declining flow to be abandoned (the phenomenon). The alternative convention (anchoring each replan's first period to the realized previous harvest, i.e., carrying the flow-history as a standing rule) suppresses the phenomenon by forcing the realized trajectory to track the plan until the constraint becomes infeasible; we implemented it as an option and report the rolling- versus shrinking-horizon and reset conventions as a robustness consideration. Second, because LP optima can be non-unique, a divergence between the projected and realized harvest vectors could in principle reflect an alternate optimum rather than genuine inconsistency. To rule this out we add an \emph{objective-gap diagnostic}: at each replan we solve the subproblem freely and again with the period-1 harvest held to the announced plan's value (within a 1\% band), and we compare the two objectives. If the re-solved objective strictly exceeds the tail-fixed objective -- or the announced value is infeasible from the realized state -- the announced plan's tail is strictly suboptimal or unimplementable, i.e., genuinely dynamically inconsistent, not merely a different member of an optimal set.

The inconsistency measure is the divergence between the open-loop plan's projected per-period harvest and the realized replanned trajectory -- that is, the changes in decisions and in projected output levels over time between what was announced and what the re-solving planner actually does. Formally, with the open-loop projection \(p = (p_1, \ldots, p_T)\) and the realized replanned trajectory \(r = (r_1, \ldots, r_T)\), we measure the per-period divergence by the symmetric relative deviation
\begin{equation}
\delta_t = \frac{|p_t - r_t|}{\max(|p_t|, |r_t|, \varepsilon)},
\label{eq:metric}
\end{equation}
which is bounded in \([0, 1]\) and robust to near-zero harvest baselines, and summarize a cell by its mean \(\bar\delta = \frac{1}{T}\sum_t \delta_t\) and by the relative change in total volume \((\sum_t r_t - \sum_t p_t)/\sum_t p_t\). A cell exhibits dynamic inconsistency when \(\bar\delta\) exceeds a small tolerance (5\%, well above LP-solver numerical noise and well below the observed magnitudes). The first-period decision is consistent by construction (the realized period-1 harvest is the open-loop period-1 decision), so the divergence is in the plan's tail, as the theory predicts. The simulator supports a rolling fixed horizon (the default, which avoids the terminal artifacts of a shrinking horizon) and a shrinking-horizon variant; all seed-dependent runs are fixed-seeded and bit-stable.

\hypertarget{experiment-design}{%
\subsection{Experiment design}\label{experiment-design}}

The experiment grid reproduces the thesis's three experimental factors: initial forest condition (landbase), harvest policy, and the interest rate. The reproduced grid spans \emph{all eighteen} of the thesis's initial forest conditions \citep[Table~5.5]{daugherty1991credibility}: the all-mature landbases (1, and 2 excluding the negatively valued CM-CE ecoclass); the area-control-derived landbases (3--8: landbase 1 or 2 after 40 or 70 years of area-control harvest at a 70- or 100-year rotation, a disequilibrium\(\to\)regulation gradient); the structured young-growth landbases (9, equal acres by age class; 10, unequal); and the randomly generated young-growth landbases (11--18, seed-fixed). Harvest policy is the thesis's six harvest-flow constraint sets \citep[Table~5.6]{daugherty1991credibility}: no harvest flow (NHF, the control); non-declining yield (NDY); bounded decline of 10\% or 20\% (\(-10\%\), \(-20\%\)); and symmetric bounded deviation of 10\% or 20\% (\(\pm10\%\), \(\pm20\%\)). The discount rate takes the thesis's four values (0\%, 2\%, 4\%, 6\%). Each of the 432 cells is a full sequential-replanning simulation under the cell's policy (equation~\eqref{eq:openloop}), scored for the occurrence and magnitude of dynamic inconsistency (equation~\eqref{eq:metric}). The grid is run by a tracked entry point and is embarrassingly parallel. (The thesis re-solved for eleven periods over its planning horizon; we run the full fifteen-period horizon, a refinement that extends the observation window without changing the mechanism.)

The thesis also used ending-period (terminal) constraints -- ending inventory at least 80\% of the regulated forest's average inventory, and final-period harvest at most 120\% of long-term sustained yield -- with its flow-constrained runs, to minimize horizon effects. We implemented and tested these; with the corrected case-study data the model does not exhibit the horizon-end liquidation they guard against, so they are non-binding here and are omitted from the reported grid (they remain available and are flagged experimental in the stack).

\hypertarget{implementation-and-reproducibility}{%
\subsection{Implementation and reproducibility}\label{implementation-and-reproducibility}}

The complete stack -- typed, provenance-recorded case-study data; the \texttt{ws3} forest model; the explicit LP builder; the sequential-replanning simulator; the experiment runner; and the validation records -- is openly available (see Software availability). The thesis's bounded-deviation sequential-flow harvest-flow constraint (equation~\eqref{eq:openloop}) required a new even-flow constraint form in \texttt{ws3} (consecutive-period anchoring with separate decrease/increase tolerances), contributed upstream and released in \texttt{ws3} v1.1.0a5; it is a reusable addition for the classic FORPLAN harvest-flow policies. The package is test-backed, the environment is locked, and the simulation records needed to regenerate every number reported below are produced by the repository's experiment entry points.
